\section{RESPUESTAS}

\subsection{I)}

{\noindent\Large\begin{align*}\displaystyle F(x, y, z) &= \sum(2, 3, 4, 5)\end{align*}}
\subsubsection{Obtener la mínima expresión booleana}
\begin{center}
\begin{tabular}{c|ccc|c}
\rowcolor[HTML]{FFCCC9} 
\cellcolor[HTML]{9AFF99}m & x & y & z & \cellcolor[HTML]{FFFFC7}F(x,y,z) \\ \hline
0                         & 0 & 0 & 0 &                                  \\
1                         & 0 & 0 & 1 &                                  \\
2                         & 0 & 1 & 0 & 1                                \\
3                         & 0 & 1 & 1 & 1                                \\
4                         & 1 & 0 & 0 & 1                                \\
5                         & 1 & 0 & 1 & 1                                \\
6                         & 1 & 1 & 0 &                                  \\
7                         & 1 & 1 & 1 &                                 
\end{tabular}
\end{center}
\noindent Los ''1'' de la función pasan, mientras que los ''0'' se los niega. Es una suma de productos de las entradas.
{\Large\begin{align*}
    F(x,y,z) &= \overline{x}\cdot y\cdot\overline{z} + \overline{x}\cdot y\cdot z + x\cdot\overline{y}\cdot\overline{z} + x\cdot\overline{y}\cdot z
\end{align*}}
\noindent Simplificación con Karnaugh:
\begin{center}
    \begin{karnaugh-map}[2][4][1][$z$][$x\qquad y$]
    \minterms{2,3,4,5}
    \autoterms[0]
    \implicant{2}{3}\implicant{4}{5}
\end{karnaugh-map}
\end{center}
\begin{itemize}[nosep] \item Primer bloque: ($m_2$ y $m_3$) cambia $z$ y hay que negar $x$.
\item Segundo bloque: ($m_4$ y $m_5$) cambia $z$ y hay que negar $y$.
\item La función simplificada queda:\end{itemize}
    {\Large\begin{align*}
    F(x,y,z) &= \overline{x}\cdot y + x\overline{y}
\end{align*}}

\subsection{II)}

{\noindent\Large\begin{align*}\displaystyle F(x, y, z) &= \sum(1, 2, 3, 5, 6, 7)\end{align*}}
\subsubsection{Obtener la mínima expresión booleana}
\begin{center}
\begin{tabular}{c|ccc|c}
\rowcolor[HTML]{FFCCC9} 
\cellcolor[HTML]{9AFF99}m & x & y & z & \cellcolor[HTML]{FFFFC7}F(x,y,z) \\ \hline
0                         & 0 & 0 & 0 &                                  \\
1                         & 0 & 0 & 1 & 1                                \\
2                         & 0 & 1 & 0 & 1                                \\
3                         & 0 & 1 & 1 & 1                                \\
4                         & 1 & 0 & 0 &                                  \\
5                         & 1 & 0 & 1 & 1                                \\
6                         & 1 & 1 & 0 & 1                                \\
7                         & 1 & 1 & 1 & 1                               
\end{tabular}
\end{center}
\noindent Los ''1'' de la función pasan, mientras que los ''0'' se los niega. Es una suma de productos de las entradas.
{\Large\begin{align*}
    F(x,y,z) &= \overline{x}\cdot\overline{y}\cdot z + \overline{x}\cdot y\cdot\overline{z} + \overline{x}\cdot y\cdot z + x\cdot\overline{y}\cdot z + x\cdot y\cdot\overline{z} + x\cdot y\cdot z
\end{align*}}
\noindent Simplificación por Karnaugh:

\begin{center}
    \begin{karnaugh-map}[2][4][1][$z$][$x\qquad y$]
        \minterms{1,2,3,5,6,7}
        \implicant{1}{5}
        \implicant{2}{7}
    \end{karnaugh-map}
\end{center}
\begin{itemize}[nosep]
    \item Primer bloque ($m_1$, $m_3$, $m_7$ y $m_5$) permanece constante $z$, el resto se sacan.
    \item Segundo bloque ($m_2$, $m_3$, $m_6$ y $m_7$) permanece constante $y$, el resto salen.
    \item La función simplificada queda:
\end{itemize}
{\Large\begin{align*}F(x,y,z) &= z + y \end{align*}}

\subsection{III)}
{\noindent\Large\begin{align*}\displaystyle F(x, y, z) &= \sum(0, 2, 4, 6)\end{align*}}
\subsubsection{Obtener la mínima expresión booleana}
\begin{center}
\begin{tabular}{c|ccc|c}
\rowcolor[HTML]{FFCCC9} 
\cellcolor[HTML]{9AFF99}$m$ & $x$ & $y$ & $z$ & \cellcolor[HTML]{FFFFC7}$F(x,y,z)$ \\ \hline
0                           & 0   & 0   & 0   & 1                                  \\
1                           & 0   & 0   & 1   &                                    \\
2                           & 0   & 1   & 0   & 1                                  \\
3                           & 0   & 1   & 1   &                                    \\
4                           & 1   & 0   & 0   & 1                                  \\
5                           & 1   & 0   & 1   &                                    \\
6                           & 1   & 1   & 0   & 1                                  \\
7                           & 1   & 1   & 1   &                                   
\end{tabular}
\end{center}

\begin{center}
    \begin{karnaugh-map}[2][4][1][$z$][$x\qquad y$]
        \minterms{0,2,4,6}
        \implicant{0}{4}
    \end{karnaugh-map}
\end{center}
\begin{itemize}\item En el bloque permanece constante $z$, por lo que se deja y hay que negarlo.\item La función simplificada queda:\end{itemize}
{\Large\begin{align*}
    F(x,y,z) &= \overline{z}
\end{align*}}

\subsection{IV)}

{\noindent\Large\begin{align*}\displaystyle F(x, y, z) = \sum(3, 4, 5, 6, 7)\end{align*}}

\subsubsection{Obtener la mínima expresión booleana}
\begin{center}
\begin{tabular}{c|ccc|c}
\rowcolor[HTML]{FFCCC9} 
\cellcolor[HTML]{9AFF99}$m$ & $x$ & $y$ & $z$ & \cellcolor[HTML]{FFFFC7}$F(x,y,z)$ \\ \hline
0                           & 0   & 0   & 0   &                                    \\
1                           & 0   & 0   & 1   &                                    \\
2                           & 0   & 1   & 0   &                                    \\
3                           & 0   & 1   & 1   & 1                                  \\
4                           & 1   & 0   & 0   & 1                                  \\
5                           & 1   & 0   & 1   & 1                                  \\
6                           & 1   & 1   & 0   & 1                                  \\
7                           & 1   & 1   & 1   & 1                                 
\end{tabular}
\end{center}

\begin{center}
    \begin{karnaugh-map}[2][4][1][$z$][$x\qquad{}y$]
        \minterms{3,4,5,6,7}
        \implicant{6}{5}\implicant{3}{7}
    \end{karnaugh-map}
\end{center}

\begin{itemize}[nosep]
    \item El primer bloque ($m_3$ y $m_7$), cambia $x$ por lo que sale y permanecen $y$ y $z$ sin negar.
    \item El segundo bloque ($m_3$, $m_7$, $m_4$ y $m_5$) permanece constante $x$ sin negar y el resto salen.
    \item La función simplificada queda:
\end{itemize}
